\chapter{市场微观结构与执行：从订单簿到事件仿真}
\label{ch:lower-microstructure-events}

\begin{itemize}
  \input{tex/generated/prerequisites/lower-microstructure-model}
  \item \textbf{学习产出：}从事件强度和 FIFO 订单簿出发，分解冲击与实施差额，求解执行和库存控制，并用可重放仿真比较策略。
  \item \textbf{研究边界：}公开逐笔成交不含完整队列状态；冲击系数、成交率和事件生成机制都需要估计，合成簿不能冒充真实成交概率。
\end{itemize}

限价单是否成交，取决于它前面的数量怎样减少。先把事件率与等待时间联系起来，再按价格和时间顺序逐笔扣减队列。


\section{事件时间不是把横轴换个名字}

令 $N_t$ 表示截至 $t$ 的事件数，历史信息为 $\mathcal F_{t^-}$。条件强度定义为
\[
 \lambda(t\mid\mathcal F_{t^-})
 =\lim_{h\downarrow0}\frac{\mathbb P(N_{t+h}-N_t=1\mid\mathcal F_{t^-})}{h}.
\]
它描述给定过去后下一瞬间的事件率，而不是“每秒平均几笔”的同义词。齐次 Poisson 过程要求独立平稳增量；等待时间 $\Delta_i$ 独立且服从参数 $\lambda$ 的指数分布。观察 $n$ 段等待时间、总暴露 $T=\sum_i\Delta_i$ 时，
\[
 \ell(\lambda)=n\log\lambda-\lambda T,
 \qquad \widehat\lambda=n/T.
\]
算例中 $n=T=4$，所以 $\widehat\lambda=1$、最大对数似然为 $-4$。这既有手算 解析参照，也由一维数值优化独立重建。

盘中 U 形季节性会让常强度模型产生系统残差。可写成
\[
 \lambda(t\mid\mathcal F_{t^-})=s(t)\widetilde\lambda(t),
\]
其中 $s(t)>0$ 是预先估计并在样本外事先确定的季节因子。若先用全日数据估计 $s(t)$ 再回测同日，会把未来活动率泄漏到早盘。

\section{Hawkes 自激与稳定边界}

指数核 Hawkes 过程写成
\[
 \lambda_t=\mu+\alpha\sum_{t_i<t}e^{-\beta(t-t_i)},
 \qquad \mu>0,\ \alpha\ge0,\ \beta>0.
\]
每个事件瞬时增加强度 $\alpha$，影响以速度 $\beta$ 衰减。分枝比 $\alpha/\beta<1$ 是平稳版本的稳定条件；它不是“拟合得好”的经验阈值。区间 $[0,T]$ 的对数似然为
\[
 \ell=\sum_{t_i\le T}\log\lambda_{t_i}
 -\mu T-\frac{\alpha}{\beta}\sum_{t_i\le T}\left(1-e^{-\beta(T-t_i)}\right).
\]
第一项奖励事件发生处的高条件强度，第二项是补偿子；漏掉补偿子会把强度无限推高。估计时还要报告窗口、初始历史、时间单位、参数约束和残差时间变换，而不能只给三个参数。

% Generated by tools/build_figures.py; edit figures/figure-specs.json instead.
\MFQFigure{tex/figures/generated/lower-ms-hawkes.pdf}{Hawkes 强度在事件到达后跳升并指数衰减；分支比接近一时，聚集持续时间与方差显著上升。}{lower-ms-hawkes}


事先确定事件时刻为 $(0.2,0.7,1.4,2.0)$，观察窗明确延伸到 $T=2.5$，参数 $(\mu,\alpha,\beta)=(0.8,0.3,1.2)$，初始激发为 $h_0=0.1$。此时上式强度还需加 $h_0e^{-\beta t}$，对数似然的补偿项还需减 $h_0(1-e^{-\beta T})/\beta$。递推实现与独立数值积分都得到对数似然 $-2.940341$；若把观察窗缩到最后事件之前，程序直接拒绝。季节调整例的基准强度为 $4/3$，变换后等待残差均值为 1。由此，观察窗、初始历史与季节暴露均不是只写在正文里的口头约定。

\begin{diagnostic}
Poisson、季节性 Poisson 与 Hawkes 是嵌套的研究假设。先检查季节性和时间尺度，再判断剩余聚集是否支持自激；不要用一个 Hawkes 名称同时解释开收盘季节性、新闻簇集和数据时间戳重复。
\end{diagnostic}

事件结构为微观结构路线 Capstone 提供可复核的队列、到达和成交假设，后续控制与仿真只能在这些边界内比较。

\section{订单流与价格变化必须在同一时间协议中}

令 $q_i\in\{-1,+1\}$ 表示第 $i$ 笔主动方向，$\Delta m_i$ 是在事先确定采样规则下的中价变化。最小联合模型为
\[
 \Delta m_i=a+bq_i+u_i,
 \qquad \mathbb E(u_i\mid q_i)=0.
\]
系数 $b$ 是给定对齐协议下的条件均值差，不自动等于结构性永久冲击。算例的 OLS 斜率为 $0.15$，透明协方差公式与带截距最小二乘一致。

真实轨采用 SEC 的公共领域订单位置摘要。它给出五个互斥位置的事件占比与 cancel-to-trade 比率\cite{sec2014orderplacement}；解析器要求占比和为 1、比率非负且所有数字有限。用 $1/(1+\text{ratio})$ 得到的只是“取消或成交”二选一近似下的隐含执行概率，加权值为 $0.044470$。合成样本 另行声明每秒 10 个终止事件的压力尺度，两者相乘才得到每秒耗尽强度；10 并非 SEC 估计。聚合摘要没有订单方向与中价变化，不能估计本节的联合回归系数。

\section{价格优先、时间优先与队列位置}

限价簿状态由订单 $(\mathrm{id},\mathrm{side},p,q,\mathrm{seq})$ 组成。被动簿至少满足：订单编号唯一、活动数量为正、最高买价低于最低卖价；同一价格按序号 FIFO。对订单 $j$，前方数量为
\[
 A_j=\sum_{i:\,p_i=p_j,\,\mathrm{seq}_i<\mathrm{seq}_j}q_i.
\]
若对手方耗尽数量在期限 $h$ 内近似 $D_h\sim\mathrm{Poisson}(\nu h)$，自己的 $Q$ 手全部成交概率是
\[
 \mathbb P(D_h\ge A_j+Q).
\]
算例 $A_j=2,Q=1,\nu h=3$，概率为 $0.576810$。这是明确假设下的队列概率，不是由公开成交快照识别出的现实概率。

% Generated by tools/build_figures.py; edit figures/figure-specs.json instead.
\MFQFigure{tex/figures/generated/lower-ms-order-book.pdf}{订单簿按价格优先、时间优先组织队列；市场单消耗可见深度并留下成交与剩余状态。}{lower-ms-order-book}


执行一笔 4 手市价卖单时，合成簿先吃掉第一张 3 手买单，再使第二张 2 手买单部分成交 1 手；随后撤去剩余 1 手，下一笔市价卖单得到未成交 1 手。这个小计算记录同时覆盖撤单、部分成交、流动性不足和状态不变量。

本单元向本路线的综合项目交付事件方向、队列状态、成交规则和可审计的时间协议\cite{hawkes1971spectra,cont2001empirical}。

\MFQTransition{时钟与标记}
事件除时刻外，还包含方向、数量与价格档。对多维买卖事件，Hawkes 稳定条件推广为核积分矩阵的谱半径小于 1；只拟合总事件数不能识别哪一侧流动性将被消耗。事件时间下相邻两条记录可能间隔极不相同，等待时间本身仍是信息。

\MFQTransition{订单流不平衡与短期价格响应}

简化 order-flow imbalance 可写为买侧新增/撤减与卖侧新增/撤减的净差。短期价格变化回归
\[
\Delta m_{t+h}=\beta\,\mathrm{OFI}_t+\varepsilon_{t+h}
\]
需要同一事件窗口、正确方向和无重叠标准误。$\beta$ 随深度和波动变化，不能从一段样本永久外推。

买卖价差补偿订单处理、库存和逆向选择。知情交易概率上升时，被动报价更容易在价格即将不利移动前成交。高成交率可能是坏消息，不能只把 fill rate 当执行质量。

\begin{diagnostic}
用聚合 K 线拟合“队列成交概率”通常不可识别：缺少订单 ID、队列前方量、撤单和到达顺序。此时应把概率标成压力假设，而不是经验估计。
\end{diagnostic}
\subsection{队列尾概率与 Hawkes 平均强度}

若每次到达恰耗尽一手，前方两手、自身一手，则全部成交要求 $N_h\geq3$。$N_h\sim\operatorname{Poisson}(3)$ 时
\[
 \mathbb P(N_h\geq3)=1-e^{-3}\left(1+3+\frac{3^2}{2}\right)\approx0.5768.
\]
若事件数量随机，就应研究标记数量之和，不能把事件笔数直接当股数。

指数 Hawkes 的平稳平均强度 $\bar\lambda$ 满足 $\bar\lambda=\mu+(\alpha/\beta)\bar\lambda$，因为单位时间的每个历史事件平均贡献核积分 $\alpha/\beta$。所以 $\bar\lambda=\mu/(1-\alpha/\beta)$。$\mu=0.8,\alpha=0.3,\beta=1.2$ 时为 $16/15$；有限窗口从空历史开始，平均强度未必立即达到平稳值。

notebook 用同价订单数量 $(3,2)$ 演示四手市价单先成交第一单 3 手、再成交第二单 1 手，并计算不同前方数量的成交概率。
